\noindent\textit{2015 manuscript.}

\section{Chemical Reactions}
Topics:
\begin{itemize}
  \item Heat of formation i.e. standard enthalpy of formation
\end{itemize}

I will follow Chapter 10, ``Chemical Reactions'', of Landau and Lifshitz (1980) \cite{LLandauELifshitz1980}.  

From Sec. 101. The condition for chemical equilibrium of Landau and Lifshitz (1980) \cite{LLandauELifshitz1980}, for a chemical reaction,
\[
\sum_i \nu_i A_i = 0 \qquad \nu_i \in \mathbb{Z}
\]
Denote thermodynamic potential $\Phi$.  

In equilibrium, say at constant temperature and pressure, 
\[
\begin{gathered}
  d\Phi\left( \frac{ \partial }{ \partial N_i} \right) = \frac{ \partial \Phi}{ \partial N_1} + \frac{ \partial \Phi}{ \partial N_2} \frac{ \partial N_2}{ \partial N_1} + \frac{ \partial \Phi}{ \partial N_3} \frac{ \partial N_3}{ \partial N_1} + \dots = 0 \text{ for }  \\
  d\Phi = \frac{ \partial \Phi}{ \partial N_1} dN_1 + \frac{ \partial \Phi}{ \partial N_2} dN_2 + \dots 
\end{gathered}
\]
where $N_1, N_2 \dots =$ number of particles of particles of various substances taking place in reaction. 

It's clear that 
\[
dN_i = \frac{ \nu_i}{ \nu_1} dN_1
\]
so 
\[
\sum_i \frac{ \partial \Phi}{ \partial N_i} \frac{ \nu_i}{ \nu_1} = 0 
\]
Let $\mu_i := \frac{ \partial \Phi}{ \partial N_i}$.  
\begin{equation}
\Longrightarrow \sum_i \nu_i \mu_i = 0 
\end{equation}

Here, I will follow Section 3.5 Chemical potential of Le Bellac, Mortessagne, Batrouni (2004) \cite{MLeBellacFMortessagneGBatrouni2004}.  

Recall that 
\[
\begin{gathered}
  G(N,\tau,p) = N\mu(\tau,p) \\
  G = \sum_j \mu_j N_j 
\end{gathered}
\]
from Eq. \ref{Eq:Gibbspotential}, and Kittel and Kroemer (1980) \cite{CKittelHKroemer1980}.  Therefore,
\[
dG = \sum_j \mu_j dN_j + N_j d\mu_j = -\sigma d\tau + Vdp + \sum_j \mu_j dN_j
\]
and so the \textbf{Gibbs-Duhem} relation is (correctly) obtained:
\begin{equation}\label{Eq:Gibbs-Duhem}
  0 = \sigma d\tau - V dp + \sum_j N_j d\mu_j
\end{equation}

\subsubsection{Conduction and Convection}\label{SubsubSec:ConductionConvection}

In $\tau d\sigma$, let's distinguish the \emph{conduction} component and \emph{convection} component.

Consider system at equilibrium, thermal equilibrium with thermostat at temperature $\tau$ \\
\qquad kept at pressure $p$ \\
\qquad able to exchange particles with reservoir at chemical potential $\mu$

Define

\[
\begin{aligned}
  & \widehat{\sigma} := \frac{ \sigma}{N } \qquad \, & \text{ \textbf{entropy per particle } } \\
  & \widehat{\epsilon} := \frac{ E}{N } \qquad \, & \text{ \textbf{energy per particle } } 
\end{aligned}
\]

Then, 
\[
\boxed{ 
\tau d\sigma = \tau d(\widehat{\sigma}N) = \tau N d\widehat{\sigma} + \tau \widehat{\sigma} dN
}
\]
where \\

$\tau N d\widehat{\sigma}$ is \emph{entropy change due to change in entropy per particle}, i.e. \emph{conduction term} and \\

$\tau \widehat{\sigma} dN$ is \emph{entropy change due to change in number of particles} i.e. \emph{ convection term}.  

So
\[
dE = Q + W + \mu dN = \tau d\sigma + 0 + \mu dN = \tau Nd\widehat{\sigma} + (\tau \widehat{\sigma} + \mu ) dN = \tau N d\widehat{\sigma} + \widehat{h} dN
\]
where
\[
\widehat{h} := \frac{H}{N} = \widehat{\epsilon} + pV \text{ is \textbf{enthalpy per particle} }
\]
since 
\[
\begin{gathered}
H = E + pV \\ 
G = E-\tau \sigma + pV = \mu N = H-\tau \sigma
\end{gathered} \Longrightarrow \tau \widehat{\sigma} + \mu = \frac{H}{N} = \widehat{h}
\]

Interpretation of $\widehat{h}dN$: If $dN$ particles are transported into a given volume by convection, energy increases by $\widehat{\epsilon}dN$.  \\
However, to return to initial volume, it's necessary to compress by $vdN =: \frac{V}{N}dN$ and add work done by pressure $pvdN$
\[
\Longrightarrow \widehat{h}dN = (\widehat{\epsilon} + pv)dN
\]
An example of this is \emph{Joule-Thomson expansion}.
\begin{correction}{enthalpy per particle}
$\hat h = H/N = \hat\epsilon + pv$ with $v = V/N$.
The display that adds $pV$ to $\hat\epsilon$ mixes an energy per particle with an energy.
See \S\ref{sec:enthalpy-per-particle}.
\end{correction}  

